\section{Direct Neural Evidence and Activated Precision Adaptation}
\label{sec:r41evidence}
The direct certificate changes both the verified continuation and the
returned policy. At each date it evaluates the trained current network
against a Bellman update using the trained next network, then stores an
action minimizing an upper interval endpoint. The terminal fit is included.
We subsequently evaluate that newly constructed policy on its own shock
tree. Its displayed values are not values of the former spline policy or
of the previously selected neural actor.

\subsection{Resolution, policy accuracy, and the economic comparison}
The design retains the four investment economies, with prices $P=1,4$ and
risk parameters $\theta=0,1$. It uses the frozen trained networks and tries
state--action resolutions $(256,128)$, $(1024,512)$, $(2048,1024)$, and,
if necessary, $(4096,2048)$. Refinement stops at the first mesh at which the
maximum certified loss at the four displayed initial states is at most
$0.02$. The first two meshes fail that query target in every economy;
the third passes in every economy. These are twelve new post-training
verification services, not twelve new training runs.

\begin{table}[htbp]\centering\small
\caption{Direct certificates for the stored neural continuations}\label{tab:r41direct}
\begin{tabular}{rrrrrr}\toprule
$P$ & $\theta$ & All-state bound & Query bound & Execution allowance & All refinements (s)\\\midrule
1 & 0 & 0.043092 & 0.016279 & 9.719 & 8.978\\
1 & 1 & 0.043334 & 0.016402 & 9.751 & 32.091\\
4 & 0 & 0.045472 & 0.017829 & 10.235 & 8.464\\
4 & 1 & 0.045813 & 0.017781 & 10.266 & 31.319\\
\bottomrule\end{tabular}
\par\smallskip\begin{minipage}{.97\textwidth}\footnotesize Notes: The query column is the maximum certified loss at $x=1/8,1/4,1/2,3/4$, not an all-state bound. The execution allowance is in units of $10^{-6}$ and covers rounded state acquisition, transition, and costs. Every final mesh has 2,048 state cells and 1,024 action cells; the clocks sum all three attempted meshes through the result-file flush and synchronization. They exclude the inherited training service and are not an end-to-end speed comparison.\end{minipage}\end{table}

The all-state loss remains between approximately $0.0431$ and $0.0459$.
The sharper displayed-state bounds, below $0.0179$, use independently
computed own-policy upper values and the direct neural lower bound on the
optimal value. They must not be extrapolated to unreported states. The
nonzero numerical-execution allowance is approximately $10^{-5}$; it is
added to the appropriate global comparison, not silently discarded.

At the accepted resolution we independently construct a conventional
spline policy and evaluate it at the same initial states. The supplement
reports each interval for neural cost minus this newly optimized spline
cost. All sixteen lie inside $[-0.001,0.001]$. Fourteen have strictly positive
lower endpoints; none has a strictly negative upper endpoint, and two
contain zero. Thus a small descriptive cost margin is established at these
states, while the conventional policy still has strictly smaller verified
cost in fourteen comparisons. Interval overlap is not used as a test of
equivalence. The margin is an explicit descriptive accounting tolerance,
not a preregistered economic indifference threshold or a population claim.

For $P=4$, gains over continuing the former $P=1$ rule have positive lower
endpoints at all eight displayed state--risk pairs, ranging from about
$0.0233$ to $0.0988$. The largest displayed-state near-optimality upper
bound is smaller than the smallest of these gain lower bounds. This
establishes that the direct neural policy's certified regret at the reported
states is smaller than the verified benefit of reoptimization there. It
does not make the price change, the benefit of reoptimization, or the
conventional competitor specific to NBO. All earlier adverse structural,
spline, and finite-law comparisons remain in the article and companion.

\subsection{When binary32 is no longer sufficient}
The earlier factorial correctly separates precision from caching, but its
adaptive paths never require binary64. A separate pressure catalogue now
crosses dimensions $2,8$, target condition numbers $1,16$, and policy targets
$10^{-10},10^{-12},10^{-14}$. The horizon is four, innovation scale is
$0.001$, and risk sensitivity is $0.1$. The inherited acceptance gates are
unchanged. Every proposed factor records its precision, input and output
hashes, acceptance decision, and target identity. A rejected binary32
proposal is charged even when a binary64 proposal subsequently succeeds.

\begin{table}[htbp]\centering\small
\caption{Precision pressure: accepted updates and rejected attempts}\label{tab:r41precision}
\begin{tabular}{lrrrrr}\toprule
Method & Certified / 12 & Accept 32 & Accept 64 & Reject 32 & Reject 64\\\midrule
adaptive cached & 12 & 116 & 42 & 42 & 0\\
fixed32 cached & 2 & 49 & 0 & 10 & 0\\
fixed64 cached & 12 & 0 & 158 & 0 & 0\\
tuned cached & 12 & 49 & 147 & 10 & 0\\
\bottomrule\end{tabular}
\par\smallskip\begin{minipage}{.97\textwidth}\footnotesize Notes: Twelve objects cross dimension and conditioning with three economic targets, $10^{-10}$, $10^{-12}$, and $10^{-14}$. Failed fixed32 services are retained. Counts for tuned precision include unsuccessful trials before restarting at the higher precision. One execution per object identifies operation of the mechanism, not a timing distribution.\end{minipage}\end{table}

Ten of the twelve adaptive services use both precisions; the other two
succeed using binary32 alone. Fixed binary32 certifies only two objects.
Adaptive and fixed binary64 each certify all twelve. The adaptive path
accepts 116 binary32 and 42 binary64 proposals and rejects 42 binary32
proposals. These observations identify actual precision escalation rather
than inferring adaptation from the controller's name. They do not establish
that adaptive complete work is smaller than the best fixed native precision:
extra attempts and verification have real costs, and this pressure catalogue
has one clock per object. A stable efficiency ranking requires a separate
matched-work study.

\subsection{Provenance and what the work comparison measures}
The source-and-result records identify the R39 reviewed article and the
partially executed R40 development attempt separately. The latter stopped
when a deserialized policy's list-valued knots were used as a numerical
array. The new driver restores both knots and actors to arrays before
repricing the old policy. It reruns the entire catalogue in a fresh directory;
it neither overwrites the failed-run outputs nor substitutes new clocks
for the 315 earlier services. The inherited protocol predates this replay,
but some R40 outcomes had already been observed. We therefore do not call
the replay a new blinded or independently prospective experiment.

All three attempted mesh clocks, failed precision attempts, interval
verification, independent queries, result serialization, and the result
file's synchronization are retained. CPU affinity and one-thread execution
are recorded; frequency is not controlled. The inherited reference-label
construction and neural training costs are not part of the new
post-training clocks and cannot be forgotten in a from-scratch comparison.
The conventional comparison is made at the accepted mesh, not at a
separately optimized common stopping target. Accordingly neither set of
new clocks is a claim of neural end-to-end acceleration.

Theorem~\ref{thm:r41dimension} in the supplement gives the explicit cost of
multidimensional verification and continuous-innovation quantization,
including the additional own-policy query tree. This theorem extends the
certificate's domain and makes its dimensional cost visible. The new
executed nonlinear example remains one-dimensional. The evidence therefore
establishes a direct neural continuation certificate and an activated native
precision mechanism, not yet a competitive high-dimensional nonlinear
performance frontier.
